\documentclass[11pt]{article}
\usepackage{fontspec}
\setmainfont{Times New Roman}
\setsansfont{Arial}
\setmonofont{Consolas}
\usepackage{amsmath,amssymb,mathtools}
\usepackage{geometry}
\usepackage{microtype}
\geometry{a4paper,margin=2.05cm}
\setlength{\parindent}{1.1em}
\setlength{\parskip}{0pt}
\newcommand{\foreign}[1]{#1}
\providecommand{\Norm}{\operatorname{N}}
\providecommand{\Sp}{\operatorname{Sp}}
\providecommand{\Gg}{\mathfrak G}
\providecommand{\Hh}{\mathfrak H}
\providecommand{\Ii}{\mathfrak I}
\providecommand{\PP}{\mathfrak P}
\providecommand{\oo}{\mathfrak o}
\providecommand{\glqq}{``}
\providecommand{\grqq}{''}
\emergencystretch=220em
\allowdisplaybreaks
\sloppy
\hbadness=10000
\vbadness=10000
\hfuzz=5pt
\vfuzz=12pt
\raggedbottom
\begin{document}

\section*{43. Диференциранье идеалов и диферента}

\begin{center}
\emph{\foreign{Journal f. d. reine u. angew. Math. 188 (1950), S. 1--21}}
\end{center}

\subsection*{Увод}

Главны теорем теорије разветвјеньа алгебраичного числового польа, како знано, говори, же диферента, идеал разветвјеньа, јест делима најманье \((e-1)\)-м степеном простого идеала, кторы в простом числу \(p\) входит в \(e\)-м степену; точно \((e-1)\)-м степеном, когда \(e\) не јест делимы с \(p\), и вышим степеном, когда \(e\) јест делимы с \(p\).

Туто јест аналог тому, же диференциалны квоциент \(f'(x)\) полинома \(f(x)\) јест делимы најманье \((e-1)\)-м степеном линеарного фактора, кторы в \(f(x)\) входит в \(e\)-м степену; точно \((e-1)\)-м степеном, когда \(e\) не јест делимы с карактеристикоју области коефициентов, и вышим степеном, когда \(e\) јест делимы с тутоју карактеристикоју.

В следујучем показујем, же ту иде о нечто вечше неж формална аналогија. Диференту можно разумети како диференциалны квоциент дефинујучего идеала числового польа \(K\) в приредјеној целочислној полиномиалној области од \(x_1,\ldots,x_n\), при чем диференциалны квоциент јест взят на месту \(x=\omega\), то јест при \(x_1=\omega_1,\ldots,x_n=\omega_n\), кде \(\omega_1,\ldots,\omega_n\) значи модулну базу система целых чисел из \(K\), главного порјадка \(\mathfrak o\). При том дефинујучи идеал \(\mathfrak M\) соотвеча целости релациј меджу \(\omega_i\), теды состојит из всех целочислних полиномов \(f(x)\), за кторе \(f(\omega)\) изчезава. Диференциалны квоциент \(\mathfrak M'[x\to\omega]\) јест дефинованы како разликовы квоциент, взят на месту \(x=\omega\).

Именно, из-за \(f(\omega)=0\), разгледајмо идеал \(\mathfrak M\) како состојаны из всех разлик \(f(x)-f(\omega)\). Далье творимо разликовы идеал
\[
        \mathfrak B=(x_1-\omega_1,\ldots,x_n-\omega_n)
\]
и разликовы квоциент
\[
        \mathfrak A=\mathfrak M:\mathfrak B,
\]
при чем творјенье квоциента значи обичны идеалны квоциент, взят в полиномиалној области од \(x\) с целыми числами из \(K\), теды с числами из \(\mathfrak o\), како коефициентами. Тогда диферента \(\mathfrak D\) дефинујетсја с
\[
        \mathfrak D=\mathfrak A[x\to\omega].
\]
Разширенье области коефициентов јест потребно, да разликовы идеал и разликовы квоциент обче были дефиноване. Иде теды о директно обобченье диференциалного квоциента полинома од једнеј неизвестној,
\[
        f'(\xi)=\frac{f(x)-f(\xi)}{x-\xi}\bigg|_{x=\xi},
\]
при чем и ту треба разширити полиномиалну област адјункцијеју \(\xi\), да разлике были дефиноване. В самој ствари идеалны диференциалны квоциент преходит в полиномиалны диференциалны квоциент, чим база \(\omega_i\) состојит из степенов једного елемента.

Дана дефиниција диференты тако директно се припојује к знаным фактам. Она але вводит неинвариантне средне чланове, поневадже идеал \(\mathfrak M\) зависи од выбраној базы. Равносилну, полно инвариантну дефиницију добивамо, когда вместо целочислној полиномиалној области вводимо колцо \(\mathfrak O\), изоморфно главному порјадку \(\mathfrak o\), кторо става изоморфно колцу класов остатков по \(\mathfrak M\), и јего област коефициентов разширјамо чрез \(\mathfrak o\) до \(\mathfrak O_{\mathfrak o}\). Разликовы идеал ту выводит се из всех разлик \((x-\xi)\), кде \(x\) и \(\xi\) соотвечајут једно другому по изоморфизму од \(\mathfrak O\) до \(\mathfrak o\); разликовы квоциент јест квоциент нулового идеала по \(\mathfrak B\); диферента взникаје из \(\mathfrak A\), когда кажды пут \(x\) заменьа се изоморфным \(\xi\).

Тута инвариантна дефиниција буде в следујучем поставјена на почеток. Тако се директно дефинује диферента произвольного комутативного колца в односу к подколцу, при чем треба само некторе условија о можности разширеньа коефициентов. То јест на примјер всегда случај при ексистенцији модулној базы, можливо после прехода к входным квоциентным колцам. Тым јест особливо дефинована и диферента произвольного порјадка числового польа, такоже диферента колца класов остатков порјадка по степену простого числа, релативна диферента и диферента при алгебраичных функцијах од неколикых неизвестных.

Совпаданье диференциалној дефиниције диференты с обычноју следује из точного структурного изследованьа Галоисового колца разширеньа, приредјеного числовому польу с помочју директного продукта. Изследованье опира се на разложенье в директну суму. В специалном случају колца класов остатков по полиному од једнеј неизвестној оно се сводит к анализи Лагрангевој интерполацијској формулы.

Точније, вводит се колцо \(\mathfrak K\), изоморфно числовому польу \(K\), кторо обхвача \(\mathfrak o\); јего област коефициентов, полье рационалных чисел, разширја се Галоисовым польем \(F\) од \(K\). Тута разширјена система \(\mathfrak K_F\) става директноју сумоју \(n\) простых компонентов, кторе соотвечајут \(n\) изоморфизмам \(K\). Ти компоненты сут разширеньа разликового квоциента и јего конјугатов; једночасно нуловы идеал \(\mathfrak K_F\) јест пресеком \(n\) идеалов, выведеных из разликового идеала и јего конјугатов. Просте компоненты имајут облик \(Fe^{(i)}\), кде \(e^{(i)}\) значет компоненты јединице. Разликовы квоциент, поневадже \(e^{(i)}\) при \(x=\xi\) преходит в јединицу, става \(\mathfrak d e^{(i)}\), кде \(\mathfrak d\) значи диференту \(\mathfrak o\). Диферента \(\mathfrak d\) состојит из всех елементов польа, кторе, помножене с \(e\), лежет в директном продукту порјадка самого с собоју. Соотвечно важет и за конјугаты.

Да из того добити дефиницију чрез комплементарны модул, треба опазити, же \(e^{(i)}\) става линеарноју формоју в бази \(x_1,\ldots,x_n\), соотвечној \(\omega\); коефициенты при \(x\) в \(e^{(i)}\) управо дају базу комплементарного модула од \(\mathfrak o\) и јего конјугатов. Из факта, же разликовы квоциент има коефициенты из \(\mathfrak o\), одмах следује, же диферента јест квоциент \(\mathfrak o\) по јего комплементарном модулу; тако добивамо Дедекиндову дефиницију. Разсмотреньа не сут ограничене на главны порјадок; оне карактеризујут диференту каждого порјадка како квоциент порјадка по јего комплементарном модулу.

Совпаданье диференциалној дефиниције с дефиницијеју чрез фундаментално равнанье следује, в случају главного порјадка, из горе поменутого факта, же идеалны диференциалны квоциент преходит в полиномиалны диференциалны квоциент, чим база состојит из степенов једного елемента, фундаменталној формы. Тым јест даны и појмовы связок тој дефиниције с Дедекиндовоју.

На основи фундаменталного равнаньа знано директно следује главны теорем поменуты на почетку. К точному определьеньу експонентов води разсмотренье диференты \(\mathfrak d[p^t]\) колца класов остатков главного порјадка по \(p^t\) с достаточно великым \(t\); достаје \(t=n\), чо јест за польа другого степена такоже нужно. Суштно јест, же диферента \(\mathfrak d\), разгледана модуло \(p^t\), точно преходит в \(\mathfrak d[p^t]\). Потом можно ограничити се на колцо класов остатков по \(\mathfrak p^t\). Оно јест изоморфно колцу класов остатков по полиному од једнеј неизвестној с коефициентами модуло \(p^t\); диференциранье тутого полинома дава, како О. Оре показал другим путем, точны преглед можливых експонентов с помочју суплементных чисел.

Наконец структурно изследованье дава још знаны теорем, же диферента надпольа од \(K\) јест продукт релативној диференты с диферентоју од \(K\). То следује из того, же иста мултипликативна властност важи за јединице разложеньа в директну суму и зато за комплементарне модулы. Соотвечне факты важет и за релативне диференты чрез преход к квоциентным колцам, такоже за алгебраичне функцијске польа с коефициентным польем карактеристики нула. При карактеристики \(p\) и при целочислних алгебраичных функцијах структурне властности сице остајут; теорија разветвјеньа але меньа се из-за појавјеньа несепарабилных поль разширеньа.

\subsection*{§1. Разширенье области коефициентов колца или директно творјенье продукта. Дефинујучи идеал}

В тутом параграфу основоју сут обча колца; комутативны закон множеньа не јест предположены. Напротив, ако противно не јест изрично зазначено, предпоклада се ексистенција јединицного елемента.

\paragraph{1. Дефиниција директного продукта.}
Колцо \(\mathfrak O\times_{\mathfrak h}\mathfrak o\), или кратше \(\mathfrak O_{\mathfrak o}\), называ се директны продукт колц \(\mathfrak O\) и \(\mathfrak o\) в односу к \(\mathfrak h\), или настало из \(\mathfrak O\) разширеньем области коефициентов \(\mathfrak h\) до \(\mathfrak o\), когда следујуче условија сут изполньене:

\begin{enumerate}
\item \(\mathfrak O_{\mathfrak o}\) садржује \(\mathfrak O\) и \(\mathfrak o\) како подколца и јест продукт тых колц.
\item Пресек \(\mathfrak O\) и \(\mathfrak o\) јест \(\mathfrak h\); при том \(\mathfrak h\) садржује јединицны елемент.
\item \(\mathfrak O\) и \(\mathfrak o\) сут елементно комутујут. Зато \(\mathfrak O_{\mathfrak o}\) состојит из билинеарных комбинациј
\[
        x_1y_1+\cdots+x_my_m,
\]
с \(x_i\in\mathfrak O\), \(y_i\in\mathfrak o\).
\item Два елементы сут једнакы точно тогда, когда они по релацијах в \(\mathfrak O\) с једнеј страны и в \(\mathfrak o\) с другој страны формално станут једнакы. Релација
\[
        \sum_i x_i y_i=0
\]
теды мора мочи бити сведена, додаваньем формално изчезајучих комбинациј облика \((yh)\delta-y(h\delta)\), на релације в обојих факторах.
\end{enumerate}

Тыми потребами равност, додаванье и множенье сут једнозначно определьене с оными факторов. Ако \(\mathfrak O\), \(\mathfrak o\) и јих пресек сут изоморфне соотвечным колцам, тогда и директны продукт јест изоморфны. Всако колцо, кторо јест продукт \(\mathfrak O\) и \(\mathfrak o\) и в ктором они елементно комутујут, јест хомоморфны образ директного продукта.

\paragraph{2. Условије ексистенције.}
Ако \(\mathfrak O\) и \(\mathfrak o\) сут колца с обчим централным подколцом \(\mathfrak h\), директны продукт не мора ексистовати. Пример: нехај \(\mathfrak h=\mathbb Z\), \(\mathfrak O=\mathfrak h[x]\) с релацијеју \(1-2x=0\), и \(\mathfrak o=\mathfrak h[y]\) с релацијеју \(1-2y=0\). Обче обхватајуче колцо с пресеком \(\mathfrak h\) не може ексистовати, бо в ньем было бы
\[
        0=(1-2x)y-x(1-2y)=y-x.
\]

Нужно и достаточно условије за ексистенцију директного продукта јест: мора ексистовати најманье једно колцо \(R\), кторо обхвача \(\mathfrak O\) и \(\mathfrak o\), в ктором јих пресек јест \(\mathfrak h\) и в ктором они елементно комутујут. Ако директны продукт ексистује, он сам јест тако колцо. Обратно, конструује се директны продукт како целост формалных билинеарных комбинациј, с равностју, додаваньем и множеньем определьеными под 1. Колцо вложеньа \(R\) тогда забезпечује условије пресека.

\paragraph{3. Дефинујучи идеал.}
Систем \(S\) елементов \(z\) из \(\mathfrak O\) называ се родны систем \(\mathfrak O\) в односу к \(\mathfrak h\), ако \(\mathfrak O=\mathfrak h[S]\). Ако \(\mathfrak h\) лежи в центру \(\mathfrak O\), тогда \(\mathfrak O\) става хомоморфны образ некомутативној полиномиалној области \(\mathfrak h\{Z\}\), кде неизвестне \(Z\) соотвечајут систему \(S\). Важи
\[
        \mathfrak O\simeq \mathfrak h\{Z\}/\mathfrak M,
\]
при чем \(\mathfrak M\) состојит из всех полиномов \(F(Z)\), за кторе \(F(z)=0\) в \(\mathfrak O\). Туты двустранны идеал \(\mathfrak M\) называ се дефинујучи идеал. Ако \(S\) состојит из једного елемента \(z\) и \(\mathfrak M\) јест главны идеал, тогда базовы полином \(G(Z)\) с \(G(z)=0\) называ се дефинујуче равнанье.

\paragraph{4. Дефинујучи идеал директного продукта.}
Нехај \(\mathfrak O\simeq\mathfrak h\{Z\}/\mathfrak M\). Поред \(\mathfrak h\{Z\}\) разгледајмо \(\mathfrak o\{Z\}\). Идеал разширеньа \(\mathfrak M_{\mathfrak o}\) состојит из линеарных комбинациј елементов \(\mathfrak M\) с коефициентами из \(\mathfrak o\). Ако директны продукт \(\mathfrak O_{\mathfrak o}\) ексистује, тогда
\[
        \mathfrak O_{\mathfrak o}\simeq \mathfrak o\{Z\}/\mathfrak M_{\mathfrak o}.
\]
Релације в директном продукту сут теды точно релације, кторе изходет из \(\mathfrak M\) разширеньем коефициентов.

\subsection*{§2. Колца с независимоју модулноју базоју. Конструкција директного продукта}

\paragraph{1. Конструкција.}
Систем \(T\) елементов \(t\) из \(\mathfrak O\) называ се \(\mathfrak h\)-модулна база \(\mathfrak O\), когда \(\mathfrak O\) состојит из всех линеарных комбинациј
\[
        h_1t_1+\cdots+h_mt_m,
        \qquad h_i\in\mathfrak h.
\]
Модулна база называ се независима, когда из \(\sum h_it_i=0\) всегда следује \(h_i=0\) за все \(i\). Ако \(\mathfrak O\) има независиму модулну базу, ктора садржује јединицны елемент, тогда директны продукт \(\mathfrak O_{\mathfrak o}\) ексистује за всако разширјајуче колцо \(\mathfrak o\) од \(\mathfrak h\), чији множествены пресек с \(\mathfrak O\) јест \(\mathfrak h\) и чији елементы комутујут с \(\mathfrak O\). Како елементы се берут все линеарны комбинације
\[
        y_1t_1+\cdots+y_mt_m,
        \qquad y_i\in\mathfrak o,
\]
с равностју коефициентов в бази \(T\). Тым јест колцо вложеньа конструовано; оно јест једночасно директны продукт. База \(T\) при том става \(\mathfrak o\)-базоју од \(\mathfrak O_{\mathfrak o}\).

\paragraph{2. Модулы разширеньа и суженьа.}
Ако \(\mathfrak B\) јест \(\mathfrak h\)-модул в \(\mathfrak O\), тогда \(\mathfrak B_{\mathfrak o}\) значи \(\mathfrak o\)-модул в \(\mathfrak O_{\mathfrak o}\), породјены ньим. Ако \(\mathfrak C\) јест \(\mathfrak o\)-модул в \(\mathfrak O_{\mathfrak o}\), тогда јего модул суженьа \([\mathfrak C,\mathfrak O]\) јест пресек с \(\mathfrak O\). Ако \(\mathfrak B\) има независиму \(\mathfrak h\)-модулну базу, ктора може бити дополньена до независимој \(\mathfrak h\)-модулној базы \(\mathfrak O\), тогда \(\mathfrak B\) јест модул суженьа својего разширеньа:
\[
        \mathfrak B=[\mathfrak B_{\mathfrak o},\mathfrak O].
\]
Доказ јест непосреднье поровнанье коефициентов в продолжаној бази.

\paragraph{3. Достаточны критериј.}
Ако ексистује разширјајуче колцо \(\mathfrak f\) од \(\mathfrak h\), тако же директне продукты \(\mathfrak O_{\mathfrak f}\) и \(\mathfrak o_{\mathfrak f}\) в односу к \(\mathfrak h\), и далеј \(\mathfrak O_{\mathfrak f}\times_{\mathfrak f}\mathfrak o_{\mathfrak f}\), ексистујут, тогда ексистује \(\mathfrak O\times_{\mathfrak h}\mathfrak o\). Бо последнье директно продуктно колцо дава колцо вложеньа, и
\[
        [\mathfrak O,\mathfrak o]\subseteq [\mathfrak O_{\mathfrak f},\mathfrak o_{\mathfrak f}]=\mathfrak f,
        \qquad [\mathfrak O,\mathfrak f]=\mathfrak h.
\]

\subsection*{§3. Диферента колца по подколцу. Диференциалны квоциент дефинујучего идеала}

Одсеј всака наступајуча колца сут комутативне. К тому колца \(\mathfrak O\) и \(\mathfrak o\) сут изоморфне, именно еквивалентне, разширеньа подколца \(\mathfrak h\), и директны продукт \(\mathfrak O_{\mathfrak o}\) в односу к \(\mathfrak h\) јест предположены.

\paragraph{1. Дефиниција диференты.}
Нехај \(x\) пробегаје все елементы \(\mathfrak O\), а \(\xi\) нека значи соотвечны јему изоморфны елемент из \(\mathfrak o\). Идеал породјены всеми разликами \(x-\xi\),
\[
        \mathfrak B=\{\ldots,x-\xi,\ldots\}
\]
называ се разликовы идеал. Разликовы квоциент јест идеалны квоциент нулового идеала по \(\mathfrak B\):
\[
        \mathfrak A=(0):\mathfrak B.
\]
Диферента \(\mathfrak d\) од \(\mathfrak o\) в односу к \(\mathfrak h\) јест
\[
        \mathfrak d=\mathfrak A[x\to\xi],
\]
а диферента \(\mathfrak D\) од \(\mathfrak O\) соотвечно \(\mathfrak A[\xi\to x]\). Обе сут изоморфно соотвечајуче идеалы.

\paragraph{2. Диференциалны квоциент дефинујучего идеала.}
Нехај \(\mathfrak M\) јест дефинујучи идеал \(\mathfrak O\) в односу к \(\mathfrak h\), соотвечно родному систему \(z\). Тогда в \(\mathfrak o[Z]\)
\[
        \mathfrak B_Z=\{\ldots,Z-\xi,\ldots\},
        \qquad
        \mathfrak C=\mathfrak M_{\mathfrak o}:\mathfrak B_Z.
\]
Теды
\[
        \mathfrak C=\{\ldots,F(Z)-F(\xi),\ldots\}:
        \{\ldots,Z-\xi,\ldots\}.
\]
Диференциалны квоциент \(\mathfrak M'\) јест \(\mathfrak C[\xi\to Z]\). Он обхвача \(\mathfrak M\) и при \(Z\to z\) односно \(Z\to\xi\) преходит в диференту \(\mathfrak O\) односно \(\mathfrak o\):
\[
        \mathfrak M'[Z\to z]=\mathfrak D,
        \qquad
        \mathfrak M'[Z\to \xi]=\mathfrak d.
\]

\paragraph{3. Случај дефинујучего равнаньа.}
Нехај \(\mathfrak O\) има дефинујуче равнанье
\[
        G(z)=0,\qquad
        G(Z)=Z^n+h_1Z^{n-1}+\cdots+h_n,
        \quad h_i\in\mathfrak h.
\]
Тогда диферента јест дефинована и јест главны идеал с базовым елементом
\[
        G'(\xi) \quad\text{односно}\quad G'(z),
\]
при чем \(G'(Z)\) јест обычны полиномиалны диференциалны квоциент. Идеалны диференциалны квоциент јест
\[
        \mathfrak M'=(G'(Z),G(Z)).
\]
В самој ствари \(1,z,\ldots,z^{n-1}\) творит независиму модулну базу. Далье
\[
        C(Z,\xi)=\frac{G(Z)-G(\xi)}{Z-\xi}
\]
јест база разликового квоциента; субституцијеју \(\xi\to Z\) добивамо \(G'(Z)\).

\subsection*{§4. Диферента директној сумы}

Нехај
\[
        \mathfrak O=R_1+\cdots+R_r,
        \qquad
        e=e_1+\cdots+e_r,
        \qquad R_iR_j=0\;(i\ne j)
\]
јест предположена како директна сума идеалов, с \(e_i\) како компонентами јединице. Далеј нехај всако \(R_i\) има независиму \(\mathfrak h_i\)-модулну базу, ктора садржује \(e_i\), при чем \(\mathfrak h_i=\mathfrak h e_i\), и нехај важи
\[
        [\mathfrak h,\mathfrak O_i]=0
\]
за соотвечне комплементарне суманды. Тогда и \(\mathfrak O\) има независиму \(\mathfrak h\)-модулну базу, творјену из баз \(R_i\), и диферента \(\mathfrak O\) в односу к \(\mathfrak h\) јест директна сума диферент \(R_i\) в односу к \(\mathfrak h_i\).

Точније: нека исто разложенье важи за \(\mathfrak o\) с компонентами \(r_i\) и идемпотентами \(\varepsilon_i\). Тогда в директном продукту
\[
        \mathfrak O_{\mathfrak o}=\sum_{i,j} R_i\varepsilon_j,
\]
и разликовы идеал се разлагаје тако, же компонента \(R_i\varepsilon_i\) садржује разликовы идеал директного продуктного творјеньа \(R_i\times r_i\). Разликовы квоциент јест директна сума соотвечных квоциентов. Зато, с \(\mathfrak A\) за разликовы квоциент и \(\mathfrak D,\mathfrak d\) за диференты,
\[
        \mathfrak D=\mathfrak D_1+\cdots+\mathfrak D_r,
        \qquad
        \mathfrak d=\mathfrak d_1+
        \cdots+\mathfrak d_r.
\]
Једине чести доказа сут поровнаньа коефициентов в независимых базах, верификација острејших релациј пресека
\[
        [\mathfrak o,\mathfrak O_i]=0,
\]
изоморфизмы
\[
        \mathfrak o\sim\mathfrak o e_i,
        \qquad
        R_i\sim R_i\varepsilon_i,
        \qquad
        \mathfrak h_i\sim \mathfrak h_i e_i,
\]
и идентификација
\[
        \mathfrak A e_i\varepsilon_i=(0):\mathfrak B e_i\varepsilon_i,
\]
теды разликового квоциента компонента. Из того следује тврдјена директно-сумна форма диференты.

\subsection*{§5. Галоисово колцо разширеньа польа. Структурне теоремы}

\paragraph{1. Основне колца.}
Нехај \(K\) јест сепарабилно полье \(n\)-того степена над основным польем \(P\). Далеј нехај \(F\) јест соотвечно Галоисово полье, кторо обхвача \(n\) конјугованых поль
\[
        K^{(1)},\ldots,K^{(n)},\qquad K=K^{(1)},
\]
и јест јих продукт. Нехај \(\mathfrak K\) значи к \(K\) изоморфно еквивалентно разширенье од \(P\), тако же \([\mathfrak K,F]=P\). Поневадже \(\mathfrak K\) има конечну независиму \(P\)-модулну базу, ексистујут директне продукты
\[
        \mathfrak K_K,
        \qquad \mathfrak K_F,
        \qquad K_F.
\]
Колцо \(\mathfrak K_F\) называ се Галоисово колцо разширеньа \(\mathfrak K\). Оно јест продукт својих \(n\) подколц \(\mathfrak K_{K^{(i)}}\).

Всакы идеал из \(\mathfrak K\) јест идеал суженьа својего разширеньа в \(\mathfrak K_K\) или \(\mathfrak K_F\); всакы идеал из \(\mathfrak K_K\) јест идеал суженьа својего разширеньа в \(\mathfrak K_F\). То следује из §2, бо над польами всакы модул има дополньајучу независиму базу.

Означимо
\[
        \mathfrak B^{(i)}=\{\ldots,x-\xi^{(i)},\ldots\}
\]
конјуговане разликове идеале и
\[
        \mathfrak A^{(i)}=(0):\mathfrak B^{(i)}
\]
конјуговане разликове квоциенты. Јих разширеньа в \(\mathfrak K_F\) будут \(\mathfrak B_F^{(i)}\) и \(\mathfrak A_F^{(i)}\).

\paragraph{2. Директно-сумно разложенье.}
Структурны теорем гласи: нуловы идеал Галоисового колца разширеньа \(\mathfrak K_F\) јест пресек разширень \(n\) конјугованых разликовых идеалов; само \(\mathfrak K_F\) јест директна сума разширень \(n\) конјугованых разликовых квоциентов:
\[
        (0)=\bigcap_{i=1}^n \mathfrak B_F^{(i)},
        \qquad
        \mathfrak K_F=\mathfrak A_F^{(1)}+\cdots+\mathfrak A_F^{(n)}.
\]
Особливо
\[
        \mathfrak K_K=\mathfrak A+\mathfrak B
\]
како директна сума, и јего нуловы идеал јест пресек соотвечного разликового идеала с разликовым квоциентом.

Доказ опира се на то, же колцо класов остатков по каждом \(\mathfrak B_F^{(i)}\) има ранг једен над \(F\), и же из сепарабилности \(\mathfrak B_F^{(i)}\) сут попарно различне и взајемно просте. Из того следује јединствено директно-сумно представјенье јединице
\[
        e=e^{(1)}+\cdots+e^{(n)},
\]
при чем
\[
        e^{(i)}\in \bigcap_{j\ne i}\mathfrak B_F^{(j)},
        \qquad
        e^{(i)}\equiv 1\pmod{\mathfrak B_F^{(i)}}.
\]
Соотвечне просте компоненты сут \(Fe^{(i)}\).

\paragraph{3. Представјенье компонентов чрез конјуговане елементы.}
Ако
\[
        x=\xi^{(1)}e^{(1)}+\cdots+\xi^{(n)}e^{(n)}
\]
јест компонентно представјенье елемента \(x\) из \(\mathfrak K\), тогда компоненты сут конјуговане. Особливо \(\xi^{(1)},\ldots,\xi^{(n)}\) сут \(n\) конјуговане вредности, кторе \(x\) пријма под изоморфизмами од \(\mathfrak K\) до \(K^{(i)}\).

Из того следује: ако \(B\) јест абсолутны модул в \(\mathfrak K\), и \(B^{(i)}\) сут јего компоненты в \(\mathfrak K_F\), тогда
\[
        B=[\mathfrak K,B^{(1)}+\cdots+B^{(n)}].
\]

\paragraph{4. Комплементарне базы.}
Всакој \(P\)-бази \(t_1,\ldots,t_n\) од \(\mathfrak K\) јест приредјена комплементарна база \(T_1,\ldots,T_n\). К ньеј изоморфне базы \(A_1^{(i)},\ldots,A_n^{(i)}\) од \(K^{(i)}\) сут определьене коефициентами компонентов јединице:
\[
        e^{(i)}=A_1^{(i)}t_1+
        \cdots+A_n^{(i)}t_n.
\]
Ако \(\alpha_1^{(i)},\ldots,\alpha_n^{(i)}\) сут конјуговане базы соотвечне \(t_j\), тогда матрице
\[
        (\alpha_j^{(i)})_{i,j}
        \quad\text{и}\quad
        (A_j^{(i)})_{i,j}
\]
сут реципрочне:
\[
        (\alpha_j^{(i)})(A_j^{(i)})^{t}=E_n.
\]
Комплементарна база комплементарној базы јест знову изходна. Ако две базы сут связане чрез \((s)=(t)P\), тогда комплементарне сут связане контрагредиентным односом
\[
        (S)=P^{-1}(T).
\]

Ако
\[
        c=C_1t_1+
        \cdots+C_nt_n
        =c_1T_1+
        \cdots+c_nT_n,
\]
тогда важи с \(\operatorname{Sp}\) како следом, т. ј. сумоју конјугатов,
\[
        C_i=\operatorname{Sp}(cT_i),
        \qquad
        c_i=\operatorname{Sp}(ct_i).
\]

\paragraph{5. Дефинујуче равнанье.}
В рукопису предвидена разработка специалного случаја дефинујучего равнаньа и Лагрангевој интерполацијској формулы остала незаполньена. Предходне теоремы дају инвариантну форму соотвечных тврдјень.

\subsection*{§6. Диферента порјадка}

Структурне теоремы за польа ныне треба заоштрити в смислу целости. Нехај \(\mathfrak h\) јест цело замкньено подколцо основного польа \(P\), чије квоциентно полье јест \(P\). Далеј нехај \(\mathfrak O\) јест \(\mathfrak h\)-порјадок в \(\mathfrak K\), теды подколцо обхвачајуче \(\mathfrak h\), чије квоциентно полье јест \(\mathfrak K\), и кторо има конечну, не нужно независиму, \(\mathfrak h\)-модулну базу. Изоморфны порјадок в \(K\) нехај буде \(\mathfrak o\), јего конјугаты \(\mathfrak o^{(i)}\); колцо породјено ними в \(F\) нехај буде \(\mathfrak g\).

Директне продукты \(\mathfrak O_{\mathfrak o}\) и \(\mathfrak O_{\mathfrak g}\) ексистујут. Диферента од \(\mathfrak O\) односно \(\mathfrak o\) в односу к \(\mathfrak h\) јест теды дефинована. Означимо \(\mathfrak B\) разликовы идеал, \(\mathfrak A=(0):\mathfrak B\) разликовы квоциент, и
\[
        \mathfrak D=\mathfrak A[\xi\to x],
        \qquad
        \mathfrak d=\mathfrak A[x\to\xi]
\]
диференты.

Структурны теорем за Галоисово колцо разширеньа порјадков гласи: разликовы идеал и разликовы квоциент порјадка преходет разширеньем в соотвечне идеале польового случаја; разликовы квоциент јест суженье разширеного квоциента. Далеј важи
\[
        \mathfrak A=\mathfrak d e^{(1)},
        \qquad
        \mathfrak A^{(i)}=\mathfrak d^{(i)}e^{(i)},
\]
и зато
\[
        \mathfrak d=[\mathfrak o,\mathfrak A^{(1)}+
        \cdots+\mathfrak A^{(n)}].
\]
Диферента \(\mathfrak d\) теды состојит из всех елементов \(x\) од \(K\), за кторе \(xe^{(1)}\) лежи в \(\mathfrak O_{\mathfrak o}\). Она јест различна од нулового идеала.

\paragraph{Комплементарны модул.}
Ако \(\mathfrak O\) има независиму \(\mathfrak h\)-базу \(t_1,\ldots,t_n\), тогда комплементарне елементы \(T_1,\ldots,T_n\) творе базу \(\mathfrak h\)-модула \(\mathfrak C\), комплементарного модула од \(\mathfrak O\). Тогда
\[
        \mathfrak d=\mathfrak o:\mathfrak C.
\]
Бо ако пишемо
\[
        e^{(1)}=A_1t_1+
        \cdots+A_nt_n,
\]
тогда \(\mathfrak d\) состојит из елементов \(d\), за кторе вси \(dA_i\) лежет в \(\mathfrak o\). В специалном случају регуларного порјадка с базоју
\[
        1,z,\ldots,z^{n-1}
\]
тако јест
\[
        \mathfrak d=(f'(z)).
\]

Еквивалентно комплементарны модул \(\mathfrak C\) состојит из всех елементов \(c\) од \(\mathfrak K\), за кторе следы
\[
        \operatorname{Sp}(co),\qquad o\in\mathfrak O,
\]
сут целе, теды в \(\mathfrak h\).

\paragraph{Додаток.}
Ако в порјадку ексистује независима \(\mathfrak h\)-модулна база, тогда разликовы квоциент \(\mathfrak A\) јест пресек всех конјугованых разликовых идеалов, различных од соотвечного разликового идеала:
\[
        \mathfrak A=[\mathfrak O_{\mathfrak o},
        \mathfrak B^{(2)}\cap\cdots\cap\mathfrak B^{(n)}].
\]
То следује из того, же тогда и разликовы идеал јест идеал суженьа својего разширеньа.

\paragraph{Фундаментално равнанье.}
Ако ексистује фундаментално равнанье, тако же после адјункције неизвестных и прехода к квоциентному колцу порјадок има базу \(1,u,\ldots,u^{n-1}\), тогда диферента јест дана полиномиалным диференциалным квоциентом дефинујучего равнаньа. Ако
\[
        G(u)=0,
\]
тогда диферента може бити прочитана из коефициентов \(G'(u)\). При целој замкньености следује потребно множенье садржаја коефициентных идеалов; тым става можлива теорија разветвјеньа чрез разсмотренье фундаменталного равнаньа модуло \(p^t\).

\subsection*{§7. Теорија разветвјеньа главного порјадка модуло \(p^t\)}

Нехај јест дано числово полье, \(\mathfrak h\) јего колцо целых чисел; или релативно полье в односу к квоциентному колцу главного порјадка по простом идеалу, при чем \(p\) значи базовы елемент тутого простого идеала. После адјункције неизвестных \(u_1,\ldots,u_n\) и главны порјадок и колцо класов остатков по \(p^t\) имајут \(\mathfrak h\)-модулну базу
\[
        1,u,\ldots,u^{n-1}.
\]
Диферента в обојих случајах јест дана с \(G'(u)\). Диферента главного порјадка преходит модуло \(p^t\) в диференту колца класов остатков по \(p^t\).

Ако
\[
        p=\mathfrak p_1^{e_1}
        \cdots\mathfrak p_r^{e_r},
\]
тогда колцо класов остатков по \(p^t\) јест директна сума компонентов, кторе соотвечајут колцам класов остатков по \(\mathfrak p_i^{te_i}\):
\[
        R=Re_1+
        \cdots+Re_r.
\]
По §4 диферента тој директној сумы јест директна сума компонентных диферент. Зато достаје разгледанье колца класов остатков по степену \(\mathfrak p^N\).

Выберимо \(\xi\) тако же
\[
        \varphi(\xi)\equiv0\pmod p,
        \qquad
        \varphi(\xi)\not\equiv0\pmod {p^2},
\]
при чем \(\varphi(x)\) јест проста функција модуло \(p\) степена \(f\). За \(t\ge2\) заменьа се \(\xi\) с \(\xi^*=\xi e_1\), кде \(e_1\) јест први идемпотент разложеньа колца класов остатков модуло \(p^t\). Тогда важи чак
\[
        \mathfrak p=(p,\varphi(\xi^*)).
\]
База модуло \(p^t\) јест дана с
\[
        \xi^{\mu}\varphi(\xi)^{\nu},
        \qquad
        \mu=0,1,\ldots,f-1,
        \quad
        \nu=0,1,\ldots,e-1.
\]
Бо из \((\varphi(\xi),p)=\mathfrak p\) и \((p)=\mathfrak p^e\) следује в колцу класов остатков
\[
        \varphi(\xi)=pM(\xi),
        \qquad M(\xi)\not\equiv0\pmod p,
\]
тако же выше степены могут сукцесивно бити изражене чрез ниже. Дефинујуче равнанье в односу к тој бази има облик
\[
        F(x)=\varphi(x)^e+pM(x),
\]
при чем \(M(x)\) модуло \(p\) не јест делимы с \(\varphi(x)\). Тым диферента колца класов остатков јест дана полиномиалным диференциалным квоциентом \(F'(\xi)\). На тутој основи определьенье експонентов своди се на резултаты О. Оре, особливо на суплементне числа и тврдјеньа Dedekind-Хенселового типа. Једночасно резултаты чрез преход к квоциентному колцу важет за релативне диференты; додаје се само изражајност диференты надпольа како продукта диференты подпольа и релативној диференты.

\begin{center}
Пријето 25 октобра 1949.
\end{center}

\end{document}

